\section*{अध्याय \ref{ch:g12:destruction-of-mountains} --- पर्वतों का विनाश}

\begin{solution}{exo:g12:destruction-of-mountains:1}
(a) \omterm{def:g12:destruction-of-mountains:denudation}{अनाच्छादन दर}; (b) \omterm{def:g12:destruction-of-mountains:load}{तल-भार}; (c) \omterm{def:g12:destruction-of-mountains:load}{घुलित भार};
(d) \omterm{def:g12:destruction-of-mountains:diagenesis}{डायजेनेसिस}; (e) \omterm{def:g12:destruction-of-mountains:rebound}{समस्थितिक प्रतिक्षेप}।
\end{solution}

\begin{solution}{exo:g12:destruction-of-mountains:2}
(1) संगुटिकाश्म। (2) बलुआ पत्थर। (3) शेल। (4) संकोणाश्म।
\end{solution}

\begin{solution}{exo:g12:destruction-of-mountains:3}
\qty{1}{m/Myr} का अर्थ है \num{1000000} वर्षों में \qty{1000}{mm}, यानी प्रति वर्ष \qty{0.001}{mm}।
(1) प्रति वर्ष \qty{0.054}{mm}। (2) प्रति वर्ष \qty{0.218}{mm}।
(3) प्रति वर्ष $\qty{10}{mm}$ का अर्थ है 1000 वर्षों में \qty{10}{m}, यानी \qty{10000}{m/Myr}।
\end{solution}

\begin{solution}{exo:g12:destruction-of-mountains:4}
ठोस भार कुल का \qty{85}{\percent} है, इसलिए कुल $60/0.85
\approx 70.6$ मिलियन टन है, और \omterm{def:g12:destruction-of-mountains:load}{घुलित भार} लगभग $70.6 - 60 \approx
10.6$ मिलियन टन प्रति वर्ष।
\end{solution}

\begin{solution}{exo:g12:destruction-of-mountains:5}
ग्रेनाइट का अपक्षय, किसी नदी द्वारा परिवहन, किसी द्रोणी में निक्षेपण,
नई परतों के नीचे दबना, सीमेंटन।
\end{solution}

\begin{solution}{exo:g12:destruction-of-mountains:6}
\qty{30}{cm/s} पर रेत-कण उठ जाता है (\qty{20}{cm/s} से ऊपर);
गाद-कण निलंबन में रहता है। \qty{5}{cm/s} पर रेत-कण बैठ जाता है
(\qty{10}{cm/s} से नीचे), पर गाद-कण अब भी निलंबन में रहता है, क्योंकि
उसे वहाँ रखने के लिए \qty{0.1}{cm/s} पर्याप्त है।
\end{solution}

\begin{solution}{exo:g12:destruction-of-mountains:7}
कुल भार $4/0.85 \approx 4.71$ मिलियन टन, यानी प्रति वर्ष $\num{4.71e9}/2830 \approx
\qty{1.66e6}{m^3}$ चट्टान। $\qty{5e10}{m^2}$ पर: प्रति वर्ष $\num{3.3e-5}$ m,
लगभग \qty{33}{m/Myr}, \qty{54}{m/Myr} की माध्यिका से कम।
\end{solution}

\begin{solution}{exo:g12:destruction-of-mountains:8}
(1) $1 - 2700/3300 \approx 0.182$। (2) लगभग \qty{182}{m}। (3) अध्याय के
\qty{163}{m} से थोड़ा अधिक: प्रावार की तुलना में भूपर्पटी जितनी हल्की,
स्तंभ उतना ही कम उठता है।
\end{solution}

\begin{solution}{exo:g12:destruction-of-mountains:9}
(1) \qty{8}{km} चट्टान। (2) $8000 \times 0.163 \approx \qty{1300}{m}$।
(3) $8000/200 = 40$ मिलियन वर्ष।
\end{solution}

\begin{solution}{exo:g12:destruction-of-mountains:10}
(1) $794/43 \approx 18$ मिलियन वर्ष। (2) सतह अपरदित मोटाई के केवल
$0.163$ गुना नीचे जाती है, इसलिए हटाई जाने वाली चट्टान $794/0.163
\approx \qty{4900}{m}$ है, जिसमें लगभग $4900/43 \approx 110$ मिलियन वर्ष लगते हैं।
(3) महाद्वीप केवल घिसते ही नहीं: पर्वत-निर्माण और
चापों की मैग्मा-क्रिया उन्हें ऊपर उठाती और नई भूपर्पटी जोड़ती रहती हैं।
\end{solution}

\begin{solution}{exo:g12:destruction-of-mountains:11}
(1) $\num{1e14}/2830 \approx \qty{3.5e10}{m^3}$, यानी लगभग
\qty{35}{km^3}। (2) $100/15.5 \approx 6.5$ वर्ष।
\end{solution}

\begin{solution}{exo:g12:destruction-of-mountains:12}
$E = kh$ और $h = h_0 e^{-fkt}$ से $E(t) = kh_0 e^{-fkt}$ मिलता है।
$t = 0$ से $t \to \infty$ तक हटाई गई मोटाई $\int_0^{\infty} kh_0
e^{-fkt}\,\mathrm{d}t = kh_0/(fk) = h_0/f$ है। सतह को $h_0$ नीचे करने के लिए
अपरदन को $h_0/f \approx 6.1\,h_0$ हटाना पड़ता है, जैसा प्रतिक्षेप वाला प्रस्ताव
कहता है: अपरदित हर मीटर सतह को केवल $f$ मीटर नीचे करता है।
\end{solution}

\begin{solution}{exo:g12:destruction-of-mountains:13}
(1) $10 \times 1000 = \qty{10000}{mm}$, \qty{10}{m}। (2) मोटाई $t$ की बर्फ़ का भार हटाना
प्रति इकाई क्षेत्रफल द्रव्यमान $\rho_{\text{बर्फ़}}\,t$ हटाने जैसा है;
स्तंभ तब तक उठता है जब तक वह उतना कम प्रावार विस्थापित न करे:
$2000 \times 900/3380 \approx \qty{530}{m}$। (3) अपने नए संतुलन के निकट आते हुए भूमि
कम से कम उठती है: उठान को चलाने वाला वह भाग है जिसकी
अभी पूर्ति नहीं हुई।
\end{solution}

\begin{solution}{exo:g12:destruction-of-mountains:14}
(1) 1900 से पहले: लगभग \qty{56}{m/Myr} (उदाहरण की तरह)।
1987--2006 में: $145/0.85 \approx 171$ मिलियन टन, $\num{1.71e11}/2830
\approx \qty{6.0e7}{m^3}$, $\qty{2.98e12}{m^2}$ पर: लगभग
\qty{20}{m/Myr}। (2) पहले वाली: बाद का भार बाँधों
और दूसरे निर्माणों द्वारा रोका जाता है। (3) नहीं: कमी प्रति वर्ष $400 - 145 = 255$ मिलियन टन
है, जिसमें से बाँध 100 से 150 की व्याख्या करते हैं, लगभग आधी।
\end{solution}

\begin{solution}{exo:g12:destruction-of-mountains:15}
(1) कम से कम \qty{6}{km}। (2) अधिकतम $310 - 252 = 58$ मिलियन वर्ष, इसलिए
कम से कम $6000/58 \approx \qty{100}{m/Myr}$। (3) यह द्रोणियों की माध्यिका
(\qty{54}{m/Myr}) और माध्य (\qty{218}{m/Myr}) के बीच है: किसी पर्वत-पट्टी के लिए
एक सामान्य दर।
\end{solution}

\begin{solution}{pb:g12:destruction-of-mountains:1}
\textbf{1.} $E = 0.209 \times 2000 \approx \qty{418}{m/Myr}$।

\textbf{2.} प्रति वर्ष \qty{0.418}{mm}।

\textbf{3.} दोनों से तेज़: यह एक खड़ी, युवा श्रेणी है, भूकंपीय रूप से सक्रिय
द्रोणियों जितनी तेज़ और उससे भी अधिक।

\textbf{4.} उसका माध्य ढाल।

\textbf{5.} $\qty{0.418e-3}{m}$ की परत, $\qty{1e6}{m^2}$ पर,
\qty{418}{m^3} है, $418 \times 2830 \approx \qty{1.18e6}{kg}$: लगभग 1200
टन।

\textbf{6.} $2000/418 \approx 4.8$ मिलियन वर्ष।

\textbf{7.} $h(t) = 2000\,e^{-0.209\,t}$ ($t$ मिलियन वर्षों में)।

\textbf{8.} $1/0.209 \approx 4.8$ मिलियन वर्ष।

\textbf{9.} $\ln 2/0.209 \approx 3.3$ मिलियन वर्ष।

\textbf{10.} $e^{-0.209 \times 11} = e^{-2.30} \approx 0.100$।

\textbf{11.} $1 - 2830/3380 \approx 0.163$।

\textbf{12.} सतह $100 \times 0.163 \approx \qty{16}{m}$ नीचे जाती है;
\omterm{def:g12:continental-crust:root}{मूल} लगभग \qty{84}{m} ऊपर उठता है।

\textbf{13.} $h' = -0.163 \times 0.209\,h \approx -0.0341\,h$, इसलिए $h(t) =
2000\,e^{-0.0341\,t}$।

\textbf{14.} $(2000 - 200)/0.163 \approx \qty{11000}{m}$ चट्टान।

\textbf{15.} वे चट्टानें जो लगभग \qty{11}{km} नीचे थीं: शिस्ट और नाइस जैसी कायांतरित
चट्टानें, और ग्रेनाइट।

\textbf{16.} $1/0.0341 \approx 29$ मिलियन वर्ष।

\textbf{17.} $29 \times \ln 10 \approx 67$ मिलियन वर्ष।

\textbf{18.} $2000\,e^{-0.0341 \times 200} \approx 2000 \times 0.0011 \approx
\qty{2}{m}$।

\textbf{19.} $2000\,e^{-0.0341 \times 250} \approx 2000 \times 0.0002 \approx
\qty{0.4}{m}$: एक समतल मैदान। \qty{2037}{m} का शिखर बताता है कि अकेला
मॉडल ठीक नहीं बैठता: तब से इस प्रदेश को फिर से ऊपर उठाया गया होगा, या वह
मॉडल की धारणा से धीमे अपरदित हुआ होगा।

\textbf{20.} माध्य ऊँचाई $29 \times \ln 2 \approx$
\textbf{20 मिलियन वर्षों} में आधी होती है।
\end{solution}
